Event-contract exchanges reproduce the oldest anomaly in betting markets. On Kalshi, longshots
lose more than half of every dollar staked on them and favourites roughly break even. The pattern
holds across sports, weather, politics, entertainment, mentions and science, at every horizon from
a month to an hour, and in liquid and illiquid markets alike. It has not shrunk as the venue has
grown by three orders of magnitude.

Who buys the longshot answers much of why the bias persists.
\begin{itemize}
  \item Small orders and small wallets buy the longshots. They do not trade at worse prices than
  large traders, but they hold most of the longshots and so bear most of the loss.
  \item Takers pay the spread and the fee on top, which together come to about a seventh of a
  longshot's stake.
  \item Liquidity suppliers take the other side at prices that just cover these costs, so the
  bias survives arbitrage without being a reliable source of profit for anyone. Wallet profits
  are extremely concentrated but do not persist.
\end{itemize}
Frictions explain why the bias is not traded away, but not why it arises. Neither insiders nor
dispersion in beliefs explains the pattern well. A single mild probability-weighting function
fits every category.

Two directions for future work follow. First, data that link identified accounts to their bet
sizes over time would separate probability weighting from risk-love and from belief dispersion,
which prices and outcomes alone cannot do. Second, the tick floor and fee schedule are design
choices. Natural experiments in either, such as a move to a finer grid near the extremes, would
show how much of the longshot premium market design can remove.
